\chapter{The Legacy Noetic Site and Its Collapse}\label{ch05:main}

This chapter is the conservative correction at the heart of the monograph. It
takes the registered legacy site---the category $\mathrm{Noetica}$ with its
three cover clauses---exactly as the conservative source states it, and proves
that two of the three clauses contribute nothing to the topology they were
intended to enrich. The ACBE clause admits, by its own reflexivity provision,
an identity arrow into its covering sieves; the RPM clause admits a split
epimorphism whose section drags an identity in by sieve closure. A sieve
containing an identity is maximal, maximal sieves belong to every Grothendieck
topology, and so each clause generates only covers that every topology already
has. The combined topology collapses to the jointly-surjective observational
coverage alone, and the modalities the two clauses were meant to seed are
trivial in the precise sense of Chapter~\ref{ch04:main}.

The result is negative but not destructive, and stating its exact perimeter is
as important as proving it. The conditional topos-existence theorem of the
conservative source survives untouched: the legacy site is still a site, its
sheaf category is still a Grothendieck topos, and nothing in this chapter
casts doubt on that. What falls is the \emph{interpretive superstructure}: the
inference that the ACBE and RPM clauses, as registered, could ever support two
distinct nontrivial modal operators. We prove the collapse in fully expanded
form (Section~\ref{ch05:sec:collapse}), diagnose it structurally---sieve
generation is a closure operator, and closure operators erase provenance
(Section~\ref{ch05:sec:closure})---and then do something the terse research
paper could not afford: we take the three most tempting repairs seriously,
formalize each, and refute each (Section~\ref{ch05:sec:repairs}). Each repair
either regenerates maximal sieves or changes nothing; the only exit is to
change the cover generators themselves, which is not a repair of the legacy
clauses but the replacement the successor chapters carry out. The chapter
closes with the exact scope of the v7.4 corrigendum: what it retracts, what it
preserves, and what discipline it imposes on all downstream text
(Section~\ref{ch05:sec:corrigendum}).

\begin{boundary}\label{ch05:bnd:status}
Status discipline. The mathematical content of this chapter consists of
elementary sieve computations whose registered status is \textsc{informal
proof present}, with the two collapse mechanisms additionally corroborated by
named executable checks in the registered finite certificate
(Section~\ref{ch05:sec:certificate}). Independent category-theory review of
these arguments remains an open obligation; nothing here is peer-reviewed,
and nothing here is promoted by appearing in this monograph. The chapter
makes no empirical, causal, safety, or product claim; the words ``ACBE'' and
``RPM'' name registered clause labels and formal structure only.
\end{boundary}

\section{The registered category $\mathrm{Noetica}$}\label{ch05:sec:noetica}

\begin{notation}\label{ch05:not:legacy}
Throughout this chapter, $\cC$ denotes the registered legacy category
$\mathrm{Noetica}$, and $J_{\mathrm{legacy}}$ denotes the coverage defined by
the conservative source's three clauses, made precise below. We write
$J_{\mathrm{js}}$ for the topology generated by the jointly-surjective
observational clause alone.
\end{notation}

The conservative source registers $\mathrm{Noetica}$ as a small concrete
category of finite sets indexed by the forty distinguished CAS indices, whose
morphisms are generated by admissibility-relation maps and closed under
pullback. Three features of this registration matter for the present chapter,
and only these three.

First, \emph{concreteness and finiteness}: the objects are finite sets, so
``jointly surjective'' has its ordinary elementwise meaning for a family of
maps into a common codomain, and the observational cover clause is a
set-level condition. Second, \emph{smallness}: $\mathrm{Noetica}$ is small in
the chosen universe, which is the hypothesis under which the background
machinery of Chapter~\ref{ch04:main}---sheafification, the sheaf topos, the
modality calculus---applies at all. Third, \emph{the registered structural
maps}: among the arrows of $\mathrm{Noetica}$ the source designates certain
families as cover-generating, and it is the exact arrows so designated, not
their intended meanings, that the topology machinery consumes. The precise
object class, the full admissibility generator list, and the index
combinatorics are registry data recorded in the frozen source artifacts and
their hashes (the governing specification is registered under SHA-256
\texttt{156A5D8C5320F36E81E3B46473BA8A038B35672DCD28839A404D96A3DB1A763A});
none of those details enter the collapse argument, which uses only the three
clause statements quoted in the next section. Where a detail of the
registration is not restated here, it is because the argument does not depend
on it, and this chapter deliberately refrains from paraphrasing registry
content it does not use.

\begin{boundary}\label{ch05:bnd:schema}
The collapse theorem is a schema-level result: it holds for \emph{every}
category and every designation of clause data satisfying the three quoted
clauses, because its proof manipulates only the clause structure (a
reflexivity provision and a section-retraction identity). In particular it
does not depend on which model of the registered axioms one takes---a point
that matters again in Chapter~\ref{ch06:main}, where a different schema-level
question (subcanonicity) is settled by exhibiting a particular model.
\end{boundary}

\section{The three cover clauses}\label{ch05:sec:clauses}

The conservative source admits a sieve as covering when it arises from any of
three clauses. We state them as registered axioms, in the notation of this
monograph, and then record the intent attached to each---keeping statement
and intent rigorously apart, because the entire content of this chapter is
that the statements do not support the intent.

\begin{axiom}[Clause (O): observational covers]\label{ch05:ax:O}
For each object $X$ of $\cC$, every family $\{f_i:U_i\to X\}_{i\in I}$ of
arrows of $\cC$ that is jointly surjective---every element of the finite set
underlying $X$ lies in the image of some $f_i$---generates a covering sieve
on $X$.
\end{axiom}

\begin{axiom}[Clause (A): ACBE descent covers]\label{ch05:ax:A}
The source designates a class of \emph{ACBE descent maps} between contexts,
and declares the sieve generated by the designated descent data at $X$ to be
covering. The designation is reflexive: at equal contexts the designated
descent map is the identity, so the generating family at $X$ contains
$\id_X$.
\end{axiom}

\begin{axiom}[Clause (R): RPM counit covers]\label{ch05:ax:R}
The RPM structure supplies, for each relevant object $X$, a pair of arrows
\[
\eta_X : X \longrightarrow R(X),
\qquad
\varepsilon_X : R(X) \longrightarrow X,
\qquad
\varepsilon_X\circ\eta_X=\id_X ,
\]
a \emph{section-retraction pair}: $\eta_X$ is a section of the split
epimorphism $\varepsilon_X$. The clause declares the sieve generated by
$\{\varepsilon_X\}$ to be covering on $X$.
\end{axiom}

\begin{remark}[Section-retraction, not adjunction]\label{ch05:rem:notadjunction}
The RPM data are registered as a section-retraction pair and expressly
\emph{not} as an adjunction. The identity $\varepsilon_X\circ\eta_X=\id_X$
holds; the other composite $\eta_X\circ\varepsilon_X$ is \emph{not} required
to be the identity of $R(X)$, and in the intended reading it typically is
not: re-embedding after realization may lose information, and prerequisite
satisfaction is one-directional. This asymmetry is a registered correction in
its own right elsewhere in the program, and the collapse argument respects
it scrupulously: the proof below uses only the one identity that is actually
postulated, $\varepsilon_X\circ\eta_X=\id_X$, and at no point assumes
$\eta_X\circ\varepsilon_X=\id_{R(X)}$, invertibility of either map, or any
adjunction triangle. The collapse is not an artifact of over-reading the RPM
structure; it follows from the weakest identity the structure carries.
\end{remark}

\begin{remark}[The recorded intent]\label{ch05:rem:intent}
The conservative source attached interpretations to the clauses: clause (O)
was to encode which families of observations jointly determine a context;
clauses (A) and (R) were to seed, downstream, two distinct modal
operators---two ``lenses'' on evidence presheaves, an ACBE-descent lens and
an RPM-resolution lens---by the mechanism of Chapter~\ref{ch04:main}: each
clause generates a topology, each topology a sheafification modality. The
acronyms are registered clause names; their operational interpretation
carries no weight in what follows and is quarantined by the program's claim
boundaries. What matters here is only the mathematical shape: (A) puts an
identity among its generators by reflexivity, and (R) puts a split
epimorphism there. Those two shapes decide everything.
\end{remark}

\subsection{Why the clauses were seductive}\label{ch05:sub:seductive}

Before refuting the construction we explain why competent
drafting produced it, because the error has a recognizable genealogy and
naming it inoculates against repetition. Both offending provisions are
natural---indeed mandatory---in \emph{adjacent} formalisms, and the legacy
source imported them into the one formalism where they are fatal.

Reflexivity first. In the theory of descent data, and everywhere equivalence
relations or groupoids organize ``being related across contexts,''
reflexivity is a harmless and structurally required axiom: an object is
related to itself, descent along the identity is trivial descent, and
omitting reflexivity would break composition. A designer specifying
``ACBE descent maps, including descent of every context to itself'' is
following that healthy instinct. The instinct fails here because a cover
clause does not feed a groupoid; it feeds a sieve, and a sieve is a
\emph{downward-closed class of arrows into a fixed target}, where an
identity member is not a harmless degenerate case but a universal solvent
(Lemma~\ref{ch05:lem:idmax}). Descent data and covering sieves are both
``families of maps with conditions,'' and the type distinction between them
is exactly what the reflexivity provision missed.

The counit second. In the RPM reading, $\varepsilon_X:R(X)\to X$ is the
realization of prerequisites---an arrow along which structured resources
map onto the target---and declaring it a cover expresses the reasonable
intent ``$X$ is covered by its resolved form.'' In formalisms where covers
are single structured morphisms with extra properties (fibrations,
effective epimorphisms taken singly, projections of resolutions), such a
declaration can carry content. Under sieve generation it cannot, whenever
the realization admits a section---and the RPM axioms themselves supply the
section $\eta_X$, since realization-after-embedding being the identity is
the whole point of the section-retraction discipline. The very axiom that
makes the RPM data well behaved (Remark~\ref{ch05:rem:notadjunction}) is
the axiom that trivializes the cover
(Lemma~\ref{ch05:lem:splitepi}). There is a genuine lesson in irony here:
the legacy construction did not fail by sloppiness in the RPM axioms---it
failed because the RPM axioms are strong enough to defeat their own
intended topological use.

\section{Two elementary facts about sieves}\label{ch05:sec:facts}

The collapse rests on two observations so short that their consequences were
easy to overlook. We isolate them, prove them, and add the one-step
generalization to isomorphisms for completeness.

\begin{lemma}[Identities force maximality]\label{ch05:lem:idmax}
Let $S$ be a sieve on an object $X$ of any category. If $\id_X\in S$, then
$S=m_X$, the maximal sieve.
\end{lemma}

\begin{proof}
Let $f:D\to X$ be any arrow with codomain $X$. Then $f=\id_X\circ f$, and $S$
is closed under precomposition with arbitrary arrows, so $f\in S$. Hence $S$
contains every arrow into $X$.
\end{proof}

\begin{lemma}[Split epimorphisms force maximality]\label{ch05:lem:splitepi}
Let $S$ be a sieve on $X$ and let $\varepsilon:E\to X$ be a split
epimorphism with section $\eta:X\to E$, i.e.\
$\varepsilon\circ\eta=\id_X$. If $\varepsilon\in S$, then $S=m_X$.
\end{lemma}

\begin{proof}
$S$ is closed under precomposition, so $\varepsilon\circ\eta\in S$; but
$\varepsilon\circ\eta=\id_X$, so $\id_X\in S$ and
Lemma~\ref{ch05:lem:idmax} applies.
\end{proof}

\begin{lemma}[Isomorphisms force maximality]\label{ch05:lem:iso}
If a sieve $S$ on $X$ contains any isomorphism $g:D\to X$, then $S=m_X$.
\end{lemma}

\begin{proof}
An isomorphism is in particular a split epimorphism, with section $g^{-1}$;
apply Lemma~\ref{ch05:lem:splitepi}. (Directly:
$g\circ g^{-1}=\id_X\in S$.)
\end{proof}

\begin{remark}\label{ch05:rem:whyblind}
Pause on \emph{why} these lemmas are inescapable. A sieve is by
definition a precomposition-closed class: the definition is what makes the
sheaf condition well posed (Definition~\ref{ch04:def:matching} quantifies
over precompositions) and what makes pullback of covers
(axiom (J2)) meaningful. Precomposition closure is therefore not an optional
convenience one might weaken to save the legacy clauses; it is constitutive
of the notion of cover in a Grothendieck topology. And precomposition closure
is exactly the rule that converts ``contains $\varepsilon_X$'' into
``contains $\id_X$'' the moment a section exists. The legacy defect is thus
not a subtle interaction of axioms but a direct consequence of the definition
of sieve meeting the definition of section.
\end{remark}

\section{The collapse theorem}\label{ch05:sec:collapse}

We first fix what ``the topology generated by the clauses'' means, so that
the collapse statement is exact.

\begin{definition}[Generated topology]\label{ch05:def:generated}
Let $G$ be a class of pairs $(X,S)$ with $S$ a sieve on $X$. The collection
of all sieves on all objects is a Grothendieck topology (the largest one),
and the intersection of any family of Grothendieck topologies is a
Grothendieck topology. Hence there is a smallest Grothendieck topology
containing $G$; we call it the topology \emph{generated} by $G$ and write
$\mathrm{Gen}(G)$. The legacy topology is
\[
J_{\mathrm{legacy}}
\;=\;
\mathrm{Gen}\bigl(G_{\mathrm{O}}\cup G_{\mathrm{A}}\cup G_{\mathrm{R}}\bigr),
\]
where $G_{\mathrm{O}}$, $G_{\mathrm{A}}$, $G_{\mathrm{R}}$ are the classes of
(object, generated sieve) pairs supplied by Axioms~\ref{ch05:ax:O},
\ref{ch05:ax:A}, \ref{ch05:ax:R} respectively, and
$J_{\mathrm{js}}=\mathrm{Gen}(G_{\mathrm{O}})$.
\end{definition}

\begin{theorem}[Legacy collapse]\label{ch05:thm:collapse}
Under the registered axioms:
\begin{enumerate}
\item every sieve in $G_{\mathrm{A}}$ is maximal;
\item every sieve in $G_{\mathrm{R}}$ is maximal;
\item consequently
$J_{\mathrm{legacy}}=J_{\mathrm{js}}$: the topology generated by all three
clauses equals the topology generated by the observational clause alone, and
the generated topology retains no ACBE or RPM provenance.
\end{enumerate}
\end{theorem}

\begin{proof}
(1) Fix an object $X$ at which clause (A) applies, and let
$S_{\mathrm{A}}(X)$ be the sieve generated by the designated ACBE descent
family at $X$. By the reflexivity provision of Axiom~\ref{ch05:ax:A}, the
designated family contains the reflexive descent map at equal contexts,
which is $\id_X$. A generated sieve contains its generators, so
$\id_X\in S_{\mathrm{A}}(X)$, and Lemma~\ref{ch05:lem:idmax} gives
$S_{\mathrm{A}}(X)=m_X$.

(2) Fix an object $X$ at which clause (R) applies, and let
$S_{\mathrm{R}}(X)$ be the sieve generated by $\{\varepsilon_X\}$. Then
$\varepsilon_X\in S_{\mathrm{R}}(X)$, and $\varepsilon_X$ is a split
epimorphism with section $\eta_X$ by the registered identity
$\varepsilon_X\circ\eta_X=\id_X$ of Axiom~\ref{ch05:ax:R}. By
Lemma~\ref{ch05:lem:splitepi}, $S_{\mathrm{R}}(X)=m_X$. Explicitly and in
full: $S_{\mathrm{R}}(X)$ is closed under precomposition, so it contains
$\varepsilon_X\circ\eta_X$; that composite \emph{is} $\id_X$; a sieve
containing $\id_X$ contains $f=\id_X\circ f$ for every $f:D\to X$; hence
$S_{\mathrm{R}}(X)$ is all of $m_X$.

(3) We show the two generated topologies coincide by showing they have the
same family of upper bounds. Let $J$ be any Grothendieck topology on $\cC$
containing $G_{\mathrm{O}}$. By axiom (J1), $J$ contains every maximal
sieve; by parts (1) and (2), $G_{\mathrm{A}}$ and $G_{\mathrm{R}}$ consist
of maximal sieves; hence $J$ contains
$G_{\mathrm{O}}\cup G_{\mathrm{A}}\cup G_{\mathrm{R}}$. Conversely any
topology containing the union contains $G_{\mathrm{O}}$. So a topology
contains $G_{\mathrm{O}}$ if and only if it contains the union, the two
intersections defining $\mathrm{Gen}$ range over the same family, and
$J_{\mathrm{legacy}}=J_{\mathrm{js}}$. For the final phrase of (3): the
equality says that the covering sieves of the legacy site are exactly the
jointly-surjective ones; no covering sieve, and hence no sheaf condition, no
sheaf, and no sheafification value, differs from what clause (O) alone
determines. Any structure downstream of the topology---the sheaf topos, the
modality, the internal logic---is therefore blind to whether clauses (A) and
(R) were ever stated.
\end{proof}

The presheaf-level restatement is given in full because it shows the
collapse operating not as an abstract maximality fact but as a concrete
gluing computation: \emph{every} presheaf satisfies the ACBE and RPM sheaf
conditions automatically, with the RPM section itself manufacturing the
amalgamation.

\begin{proposition}[Clause-wise gluing is automatic]\label{ch05:prop:autoglue}
Let $F$ be any presheaf on $\cC$.
\begin{enumerate}
\item For the ACBE generating family at $X$ (which contains $\id_X$), the
gluing map of $F$ is a bijection.
\item For the RPM generating family $\{\varepsilon_X\}$ at $X$, the gluing
map of $F$ is a bijection; explicitly, a matching datum
$x\in F(R(X))$ compatible over $\varepsilon_X$ has the unique amalgamation
$y=x|_{\eta_X}\in F(X)$.
\end{enumerate}
\end{proposition}

\begin{proof}
By the standard reduction (Chapter~\ref{ch04:main}), matching families for a
generating family biject with matching families for the sieve it generates,
so part (1) is Example~\ref{ch04:ex:maximal} applied to
$S_{\mathrm{A}}(X)=m_X$: a matching family on the maximal sieve is
determined by its value at $\id_X$, and gluing is the corresponding
bijection.

Part (2) can be quoted the same way, but we also carry out the computation
directly, because it exhibits the mechanism. A matching datum for
$\{\varepsilon_X\}$ is an element $x\in F(R(X))$ satisfying the
compatibility condition: whenever $h,h':E\to R(X)$ obey
$\varepsilon_X\circ h=\varepsilon_X\circ h'$, then $x|_h=x|_{h'}$. Set
$y:=x|_{\eta_X}\in F(X)$. We claim $y|_{\varepsilon_X}=x$. By contravariant
functoriality, $y|_{\varepsilon_X}=x|_{\eta_X\circ\varepsilon_X}$
(restriction along $\eta_X$ then along $\varepsilon_X$ is restriction along
the composite $\eta_X\circ\varepsilon_X:R(X)\to R(X)$). Now apply the
compatibility condition with $h=\eta_X\circ\varepsilon_X$ and
$h'=\id_{R(X)}$: indeed
$\varepsilon_X\circ(\eta_X\circ\varepsilon_X)
=(\varepsilon_X\circ\eta_X)\circ\varepsilon_X=\varepsilon_X
=\varepsilon_X\circ\id_{R(X)}$,
so $x|_{\eta_X\circ\varepsilon_X}=x|_{\id_{R(X)}}=x$. Hence
$y|_{\varepsilon_X}=x$: $y$ is an amalgamation. For uniqueness, suppose
$y'\in F(X)$ also has $y'|_{\varepsilon_X}=x$. Then
\[
y' \;=\; y'|_{\id_X}
\;=\; y'|_{\varepsilon_X\circ\eta_X}
\;=\; \bigl(y'|_{\varepsilon_X}\bigr)\big|_{\eta_X}
\;=\; x|_{\eta_X} \;=\; y .
\]
So gluing along the RPM family is a bijection for every presheaf $F$: the
sheaf condition it imposes is empty. Note where the two halves came from:
\emph{existence} of the amalgamation used the compatibility condition once,
at the pair $(\eta_X\circ\varepsilon_X,\ \id_{R(X)})$; \emph{uniqueness}
used only the section identity. The very map $\eta_X$ that was meant to
express one-directional prerequisite structure is the map that trivializes
the cover.
\end{proof}

\begin{corollary}[Triviality of the intended modalities]\label{ch05:cor:trivial}
Let $J_{\mathrm{A}}^{\mathrm{leg}}=\mathrm{Gen}(G_{\mathrm{A}})$ and
$J_{\mathrm{R}}^{\mathrm{leg}}=\mathrm{Gen}(G_{\mathrm{R}})$ be the
topologies generated by clause (A) alone and clause (R) alone. Then every
covering sieve of either topology is maximal; every presheaf on $\cC$ is a
sheaf for both; both sheafification modalities are trivial (naturally
isomorphic to the identity of $\PSh(\cC)$); and the two modalities are
naturally isomorphic to each other. Moreover, relative to the full legacy
site, adjoining clauses (A) and (R) to clause (O) changes neither the sheaf
category, nor the sheafification functor, nor the resulting topos:
\[
\Sh(\cC,J_{\mathrm{legacy}})\;=\;\Sh(\cC,J_{\mathrm{js}})
\quad\text{with equal, not merely equivalent, subcategories of }\PSh(\cC).
\]
\end{corollary}

\begin{proof}
The generators of $J_{\mathrm{A}}^{\mathrm{leg}}$ and
$J_{\mathrm{R}}^{\mathrm{leg}}$ are maximal sieves
(Theorem~\ref{ch05:thm:collapse}(1),(2)), and the collection consisting of
exactly the maximal sieves on each object is already a Grothendieck topology
(axiom (J1) holds by construction; pullback of a maximal sieve along $g$
contains $\id$, hence is maximal, giving (J2); and (J3) adds nothing new,
since if $T$ has $f^{*}T$ maximal for all $f$ in a maximal sieve, take
$f=\id$ to get $T$ maximal). Being the smallest Grothendieck topology on
$\cC$, it is therefore the generated one.
Proposition~\ref{ch04:prop:trivial} then gives triviality of
both modalities, and trivial modalities are naturally isomorphic to the
identity, hence to each other. The final statement is
Theorem~\ref{ch05:thm:collapse}(3): equal topologies have identical sheaf
conditions, hence identical sheaf subcategories, and sheafification is
determined up to canonical isomorphism by its subcategory
(Proposition~\ref{ch04:prop:uniqueness}).
\end{proof}

\begin{remark}[What the collapse is not]\label{ch05:rem:whatnot}
Three misreadings should be blocked immediately. The collapse is \emph{not} a
contradiction in the legacy axioms: the axioms are consistent, the site is a
site, and the topos exists. It is \emph{not} a defect of topos theory or an
exotic boundary case: Lemmas~\ref{ch05:lem:idmax}--\ref{ch05:lem:splitepi}
are among the first exercises one can state about sieves, and the theorem is
their immediate application. And it is \emph{not} a statement that ACBE or
RPM structure is meaningless: the descent maps and the section-retraction
pair remain perfectly good arrows and identities of $\cC$, available to any
construction that consumes arrows. What collapses is specifically the
attempt to make the \emph{topology} remember them, because they were fed to
it in a form (containing an identity; containing a split epi) that
generation immediately saturates. The failure is a type error, in the
following sense: cover generation is the wrong interface for provenance, and
the next section makes that precise.
\end{remark}

\section{Generation is a closure operator, and closure erases
provenance}\label{ch05:sec:closure}

The collapse theorem shows two particular clauses dissolving into the
topology. The deeper point is that \emph{any} clause information beyond the
bare arrow families must dissolve, because of what kind of operation
topology generation is. This section formalizes that claim; it is the
mathematical justification for the successor's architectural decision to
carry provenance in parallel records outside the topology, rather than
trying again with cleverer generators.

Consider the complete lattice $\mathcal{L}$ whose elements are classes $G$ of
(object, sieve) pairs on $\cC$, ordered by inclusion. Grothendieck topologies
are particular elements of $\mathcal{L}$, and, as noted in
Definition~\ref{ch05:def:generated}, they are closed under arbitrary
intersection and include the top element. The map
$\mathrm{Gen}:\mathcal{L}\to\mathcal{L}$ sending $G$ to the smallest topology
containing it is therefore the closure operator associated with a closure
system, and it enjoys the three defining properties:

\begin{proposition}[Generation is a closure operator]\label{ch05:prop:closure}
For all $G,G'\in\mathcal{L}$:
\begin{enumerate}
\item (Extensive) $G\subseteq\mathrm{Gen}(G)$;
\item (Monotone) $G\subseteq G'$ implies
$\mathrm{Gen}(G)\subseteq\mathrm{Gen}(G')$;
\item (Idempotent) $\mathrm{Gen}(\mathrm{Gen}(G))=\mathrm{Gen}(G)$;
\end{enumerate}
and the fixed points of $\mathrm{Gen}$ are exactly the Grothendieck
topologies.
\end{proposition}

\begin{proof}
All four statements are the standard facts about the closure operator of a
closure system: $\mathrm{Gen}(G)$ is the intersection of all topologies
containing $G$; that intersection contains $G$ (extensivity); a larger $G$
is contained in fewer topologies, so its intersection is larger
(monotonicity); $\mathrm{Gen}(G)$ is itself a topology, so the topologies
containing it are the topologies containing $G$ that contain it, and the
intersection is unchanged (idempotency); and $G$ is a fixed point iff $G$ is
a topology, since topologies are exactly the elements containing themselves
as their own closure.
\end{proof}

Now observe what the \emph{input} to $\mathrm{Gen}$ is. A clause system, as
registered, is richer than an element of $\mathcal{L}$: it is a family of
\emph{labeled} designations---this family is admitted \emph{as an
observational cover}, this one \emph{as ACBE descent}, this one \emph{as the
RPM counit}, each with its authority, its assumptions, its intended
downstream use. Formally, a clause system is a set $\Lambda$ of labels
together with a function $u:\Lambda\to\mathcal{L}$ assigning to each label
the (object, sieve) pairs it contributes. The topology the site carries is
\[
J \;=\; \mathrm{Gen}\Bigl(\ \bigcup_{\lambda\in\Lambda}u(\lambda)\ \Bigr).
\]
Two erasures happen in this formula, in order. The union forgets the
\emph{partition} of the generators by label: two clause systems with the same
underlying pairs but different labelings have the same union. Then
$\mathrm{Gen}$ forgets the \emph{generators themselves}: it returns a fixed
point, and a fixed point carries no record of which of its elements were
inputs and which arose from saturation. Both erasures are structural, not
accidental:

\begin{finding}[Provenance erasure]\label{ch05:fnd:erasure}
The passage from a labeled clause system to its Grothendieck topology
factors as
\[
\text{labeled clauses}
\xrightarrow{\ \text{forget labels}\ }
\text{generator classes in }\mathcal{L}
\xrightarrow{\ \mathrm{Gen}\ }
\text{topologies},
\]
and each factor is non-injective in the relevant way: the first identifies
all labelings of the same generators; the second identifies all generator
classes with the same closure. Consequently no construction that consumes
only the topology---sheaves, sheafification, the modality, the internal
logic, the topos---can distinguish two clause systems whose generators have
equal closure, and \emph{a fortiori} none can recover which clause admitted
a given cover. A Grothendieck topology is constitutionally incapable of
remembering why a family was admitted.
\end{finding}

\begin{example}\label{ch05:ex:noninjective}
Non-injectivity of $\mathrm{Gen}$ is easy to witness even without the
collapse: adding to any generator class all maximal sieves, or any sieve
already in its closure, changes nothing. The legacy site is the pointed
instance: $G_{\mathrm{O}}$ and
$G_{\mathrm{O}}\cup G_{\mathrm{A}}\cup G_{\mathrm{R}}$ are different
generator classes---materially different, in that the latter records the
program's two interpretive clauses---with identical closure
(Theorem~\ref{ch05:thm:collapse}). The interpretive investment in clauses
(A) and (R) was placed entirely in the part of the input that the closure
operator is guaranteed to forget.
\end{example}

\begin{remark}[The design consequence, stated once]\label{ch05:rem:design}
Finding~\ref{ch05:fnd:erasure} converts the collapse from an accident into a
specification. If cover provenance---which authority admitted a family,
under which assumptions, with which review state---is to survive, it must be
stored in a structure that is \emph{not} downstream of $\mathrm{Gen}$: a
parallel record bound to the exact member arrows, versioned and
hash-chained, consumed by governance code while the topology consumes the
bare family. That is precisely the marked-cover architecture the successor
adopts, developed in the marked-covers chapter of this monograph; nothing in
the present chapter depends on it, but the reader should carry forward the
exact division of labor it implements: \emph{the topology is for
compatibility semantics; the mark record is for provenance; neither can do
the other's job}.
\end{remark}

\section{Three tempting repairs, formalized and refuted}
\label{ch05:sec:repairs}

Faced with Theorem~\ref{ch05:thm:collapse}, three repairs suggest themselves
so naturally that each was, at one time or another, informally proposed
within the program. Each keeps the legacy site ``as intended'' while
adjusting the offending detail. We formalize all three and show that each
either still yields maximal sieves or changes nothing. The section is
deliberately thorough: the refutations, not the collapse itself, are what
justify the successor's decision to abandon the legacy generators rather
than continue repairing them.

\subsection{Repair 1: relabel the generators}\label{ch05:sub:repair1}

\emph{Proposal.} Keep the category and the designated arrows, but rename or
re-partition the clauses: introduce new clause labels (say
ACBE-prime and RPM-prime), or redistribute the same designated arrows among
differently named clauses, or designate the same arrows ``with different
intended semantics,'' expecting the resulting topologies to differ and the
modalities to separate.

\begin{proposition}[Generation is label-blind]\label{ch05:prop:repair1}
Let $(\Lambda,u)$ and $(\Lambda',u')$ be two labeled clause systems with the
same underlying generators:
$\bigcup_{\lambda}u(\lambda)=\bigcup_{\lambda'}u'(\lambda')$. Then they
generate the same topology, the same sheaf category, and canonically
isomorphic modalities. In particular, any relabeling of the legacy clauses
that preserves the designated arrows yields exactly
$J_{\mathrm{js}}$ again; and any relabeled clause whose designated family
still contains the reflexive descent map, or still contains a split
epimorphism onto its target, still generates only maximal sieves.
\end{proposition}

\begin{proof}
The first sentence is immediate from the factorization in
Finding~\ref{ch05:fnd:erasure}: the topology depends on the clause system
only through the union of its generators. The second sentence is
Theorem~\ref{ch05:thm:collapse} applied verbatim: the collapse proof used no
property of the labels, only membership of $\id_X$ (via reflexivity) or of a
split epimorphism (via the section identity) in the generating family, and
those memberships are properties of the arrows, which the relabeling
preserves.
\end{proof}

\begin{remark}\label{ch05:rem:repair1}
The refutation is short precisely because the proposal misidentifies where
the defect lives. The names ``ACBE'' and ``RPM'' never reached the
mathematics; only the arrow families did, and they are unchanged. A variant
of the proposal designates \emph{different} arrows under the old names---but
then either the new families still satisfy the reflexivity and
section-retraction provisions that the registered axioms impose on those
clauses (collapse again, by the same two lemmas), or the provisions have
been dropped, in which case the axioms have been changed and the proposal is
no longer a relabeling but a new site, to be analyzed on its own terms. In
either reading, relabeling per se accomplishes nothing.
\end{remark}

\subsection{Repair 2: restrict to sub-sieves}\label{ch05:sub:repair2}

\emph{Proposal.} The full generated sieves are maximal only because closure
drags in the identity; so declare covering, instead, a proper
\emph{sub-collection} of the generated sieve---the RPM arrows ``without''
the identity they generate, or the ACBE family ``without'' its reflexive
member---and expect nonmaximal covers.

\begin{proposition}[Sub-sieve restriction is unstable or vacuous]
\label{ch05:prop:repair2}
Fix $X$ and the RPM data of Axiom~\ref{ch05:ax:R}.
\begin{enumerate}
\item Any sieve on $X$ containing $\varepsilon_X$ is maximal; hence there is
no proper sub-\emph{sieve} of $m_X$ containing $\varepsilon_X$, and any
covering sieve that retains the clause's designated arrow is maximal.
\item The precomposition closure of $\{\varepsilon_X\}$ with $\id_X$ deleted,
\[
S^{-} \;=\;
\{\varepsilon_X\circ h \mid h \text{ an arrow into } R(X)\}
\setminus\{\id_X\},
\]
is not a sieve: it contains $\varepsilon_X$ but not
$\varepsilon_X\circ\eta_X$. Declaring it covering is not admissible data
for a Grothendieck topology, and applying $\mathrm{Gen}$ to it restores
$\id_X$ and yields $m_X$ again.
\item If instead the clause is weakened to designate a family omitting
$\varepsilon_X$ (respectively, omitting the reflexive ACBE map at every
context where it appears), then at every object where the omission leaves
the clause with \emph{no designation at all}, or leaves a family whose
generated sieve already lies in $J_{\mathrm{js}}$, the clause contributes
no covering sieve beyond $J_{\mathrm{js}}$: the modification changes
nothing. The remaining reading---retaining the clause but declaring the
\emph{empty family} covering at a nonempty $X$---is not a weakening at all:
the empty sieve is then covering, the hypothesis of (J3) is vacuous over
it, so every sieve on $X$ becomes covering and every sheaf $F$ is forced to
$|F(X)|=1$.
\end{enumerate}
Consequently every variant of the proposal either regenerates maximal
sieves, trivializes the topology at $X$, or reduces the clause to a vacuous
appendix of clause (O).
\end{proposition}

\begin{proof}
(1) is Lemma~\ref{ch05:lem:splitepi}. (2): membership of $\varepsilon_X$ and
closure under precomposition would force
$\varepsilon_X\circ\eta_X=\id_X\in S^{-}$, contradicting the deletion; so
$S^{-}$ fails the sieve axiom at exactly the composite the repair tried to
excise. A Grothendieck topology's covers are sieves by definition, so
$S^{-}$ cannot be declared covering as it stands; the two available
completions are to close it (restoring $\id_X$, hence maximality by
Lemma~\ref{ch05:lem:idmax}) or to shrink it until it is
precomposition-closed without $\varepsilon_X$---which is case (3). For (3):
a clause that no longer designates $\varepsilon_X$ (or the reflexive map) no
longer contains the arrows that defined it; whatever family it now
designates either generates sieves already covering in $J_{\mathrm{js}}$
(no change to the topology, by idempotency of $\mathrm{Gen}$), or generates
genuinely new covers---but then the clause is a \emph{new} cover axiom about
\emph{other} arrows, whose justification, pullback-stability, and
consistency with the registered axioms must be established from scratch.
Within the registered clause data, no such family is specified; the
registered content of clauses (A) and (R) beyond their identity and counit
members is exactly the designation recorded in the source, and the source
attaches the reflexivity and counit provisions to it. Deleting those
provisions deletes the clause's registered content. The one remaining
reading is disposed of directly: if the clause is retained but designates
the empty family at a nonempty $X$, the empty sieve is covering, so the
hypothesis of (J3) is vacuous over it and every sieve on $X$ becomes
covering, while the sheaf condition against the empty sieve forces
$|F(X)|=1$ for every sheaf $F$.
\end{proof}

\begin{remark}[The pretopology variant]\label{ch05:rem:pretop}
A more sophisticated version of the proposal keeps families, not sieves, as
the covering data---``use $\{\varepsilon_X\}$ as a pretopology cover and
never close it.'' This changes nothing either, and the reason is structural:
the \emph{sheaf condition} for a family is equivalent to the sheaf condition
for the sieve it generates (matching families biject, as used in
Proposition~\ref{ch05:prop:autoglue}), so even a formalism that never
materializes the closed sieve imposes exactly the closed sieve's gluing
requirement---which, by Proposition~\ref{ch05:prop:autoglue}(2), every
presheaf satisfies automatically. The collapse is visible at the level of
what the covers \emph{demand}, not merely at the level of how they are
stored: a cover containing a split epimorphism demands nothing, however the
bookkeeping is arranged. There is no formalism-shopping exit.
\end{remark}

\subsection{Repair 3: add new objects without changing covers}
\label{ch05:sub:repair3}

\emph{Proposal.} Embed $\cC$ into a larger category $\cC'$---adjoin
auxiliary context objects, resolution stages, or interpolants---while
keeping the three clauses as they are, expecting the richer ambient
structure to make the ACBE/RPM sieves nonmaximal (``there will be more
arrows, so the old sieves will no longer be everything'') or otherwise to
reactivate the lost distinction.

\begin{proposition}[Object adjunction preserves the collapse]
\label{ch05:prop:repair3}
Let $\iota:\cC\to\cC'$ be a full embedding into any category $\cC'$, and
suppose the clause data are transported along $\iota$ unchanged: the
designated ACBE families (with reflexivity) and RPM data (with
$\varepsilon_X\circ\eta_X=\id_X$) generate covers on the images
$\iota(X)$, and no new cover designations are made. Then:
\begin{enumerate}
\item the sieve generated in $\cC'$ by the transported ACBE family at
$\iota(X)$ is the maximal sieve on $\iota(X)$ in $\cC'$;
\item likewise for the transported RPM family;
\item objects of $\cC'$ outside the image receive no covering sieves from
the clauses at all (only what (J1) and any transported (O)-covers give
them);
\item hence the generated topology on $\cC'$ bears the same relation to its
observational part as before: the (A)- and (R)-clauses contribute only
maximal sieves, and the collapse statement holds verbatim with $\cC$
replaced by $\cC'$.
\end{enumerate}
\end{proposition}

\begin{proof}
(1) and (2): the two collapse lemmas are absolute, not relative. The proof
that a sieve containing $\id_{\iota(X)}$ contains every arrow into
$\iota(X)$ quantifies over the arrows of the \emph{ambient} category and
uses only precomposition closure; more arrows in $\cC'$ means more members
forced into the sieve, never fewer. Concretely: functors preserve
identities and composites, so the transported reflexive member is
$\id_{\iota(X)}$ and the transported RPM identity reads
$\iota(\varepsilon_X)\circ\iota(\eta_X)=\id_{\iota(X)}$; Lemmas
\ref{ch05:lem:idmax} and \ref{ch05:lem:splitepi} then apply in $\cC'$
exactly as in $\cC$. The intuition behind the proposal---that enlarging the
category ``dilutes'' a sieve---has the logic backwards: maximality of a
sieve generated by an identity is stable under arbitrary enlargement of the
ambient arrow supply, because the identity swallows whatever arrows the
enlargement provides. (3) is immediate: a cover designation at $\iota(X)$
is data at $\iota(X)$; new objects received no designations by hypothesis.
(4) restates Theorem~\ref{ch05:thm:collapse}(3) in $\cC'$, whose proof
repeats word for word.
\end{proof}

\begin{remark}\label{ch05:rem:repair3}
The proposal has a second reading: adjoin new objects \emph{and} give them
new covers. That, of course, can produce nontrivial modalities---it is
exactly what the successor does, adjoining a finite context factor whose
partition covers are proper by construction. But it falls outside the
repair's own terms (``without changing covers''), and honesty requires
saying so plainly: once new cover generators are introduced, the legacy
clauses are not being repaired; they are being retired as topology inputs,
surviving only as labels in the parallel provenance layer. The refutation
of Repair 3 is thus also the derivation of the successor's design from
first principles: since object adjunction alone changes nothing and
generator surgery within the clauses is unstable or vacuous
(Propositions~\ref{ch05:prop:repair2}--\ref{ch05:prop:repair3}), any
successor with genuinely distinct modalities must introduce new proper
generators somewhere, and the collapse lemmas dictate the design rule they
must obey.
\end{remark}

\begin{table}[ht]
\centering
{\small
\begin{tabular}{@{}>{\raggedright\arraybackslash}p{0.21\textwidth}
>{\raggedright\arraybackslash}p{0.15\textwidth}
>{\raggedright\arraybackslash}p{0.27\textwidth}
>{\raggedright\arraybackslash}p{0.20\textwidth}@{}}
\toprule
Repair & Formalization & Mechanism of failure & Verdict \\
\midrule
Relabel generators & Prop.~\ref{ch05:prop:repair1} & generation is
label-blind & changes nothing \\
Restrict to sub-sieves & Prop.~\ref{ch05:prop:repair2} & closure restores
$\id$; else clause vacuous & maximal or nothing \\
Keep bare families & Rem.~\ref{ch05:rem:pretop} & family condition equals
sieve condition & changes nothing \\
Add objects, same covers & Prop.~\ref{ch05:prop:repair3} & maximality is
absolute under embedding & changes nothing \\
New objects, new covers & Rem.~\ref{ch05:rem:repair3} & --- & not a repair;
successor route \\
\bottomrule
\end{tabular}}
\caption{The tempting repairs and their fates. Every route that keeps the
registered generators yields maximal sieves or leaves the topology
untouched; the only escape changes the generators, which is the successor
construction, not a repair.}
\label{ch05:tab:repairs}
\end{table}

\begin{invariant}[Design rule for nontrivial cover generators]
\label{ch05:inv:designrule}
A generating family intended to produce a nonmaximal covering sieve on $X$
must contain no identity of $X$, no isomorphism onto $X$, and no split
epimorphism onto $X$; on a poset, equivalently, every generating member must
be a proper subobject of $X$. This invariant is extracted directly from
Lemmas~\ref{ch05:lem:idmax}--\ref{ch05:lem:iso} and governs the successor's
partition coverages, whose global cover pieces are proper subcontexts by
construction and which therefore escape the collapse at exactly this point.
\end{invariant}

\section{What the v7.4 corrigendum retracts, and what it preserves}
\label{ch05:sec:corrigendum}

The v7.4 corrigendum is the registered document through which the program
absorbed Theorem~\ref{ch05:thm:collapse}. Because corrigenda are easily
over-read---both by critics, as a concession of more than was conceded, and
by proponents, as conceding less---this section states its exact scope in
both directions.

\subsection{Retracted}

First, the corrigendum retracts the inference that the legacy ACBE and RPM
clauses generate distinct nontrivial modalities. By
Corollary~\ref{ch05:cor:trivial} the clause-generated topologies have only
maximal covers, both associated modalities are naturally isomorphic to the
identity, and every statement in earlier informal material that spoke of
``the ACBE lens'' or ``the RPM lens'' as operators acting nontrivially on
legacy evidence presheaves is withdrawn without residue. There is no weaker
salvage---no ``partially distinct,'' no ``distinct on a subclass of
presheaves'': a trivial modality fixes every presheaf up to canonical
isomorphism, and two trivial modalities are isomorphic to each other.

Second, and consequentially, the corrigendum retracts the inference that the
legacy topology retains ACBE or RPM provenance in any form. This is
Finding~\ref{ch05:fnd:erasure} applied to the registered clauses: the
topology equals $J_{\mathrm{js}}$, and nothing downstream of the topology
can testify that clauses (A) and (R) were ever part of the site
description.

Third---developed in full in Chapter~\ref{ch06:main} and only registered
here---the corrigendum retracts any silent treatment of representable
presheaves as sheaves for the legacy topology. Subcanonicity of
$(\cC,J_{\mathrm{legacy}})$ is not derivable from the registered axioms, so
every use of a bare $yX$ as an object of the sheaf topos is replaced by its
associated sheaf.

\subsection{Preserved}

First, \emph{conditional topos existence stands}. The legacy site is a
genuine site: $J_{\mathrm{legacy}}$ is a Grothendieck topology (a collapsed
one, but the collapse is an equality of topologies, not a failure of the
axioms (J1)--(J3)), and for a small site the sheaf category is a
Grothendieck topos by the standard background theorem
(Theorem~\ref{ch04:thm:sheafification}; \cite{maclane,johnstone}). The
conditional statement---\emph{if} the registered site data are well formed
and small, \emph{then} $\Sh(\cC,J_{\mathrm{legacy}})$ is a Grothendieck
topos---was true before the collapse was noticed and is true after. The
corrigendum takes care to permit this statement while blocking the modal
inference; the two have entirely different evidentiary bases.

Second, the \emph{observational coverage survives with its content intact}.
$J_{\mathrm{legacy}}=J_{\mathrm{js}}$ means the legacy sheaf topos is
exactly the topos of jointly-surjective descent: everything the legacy
construction validly proved about compatibility of observations over
jointly-surjective families remains proved, and the legacy topos remains the
reference object for the successor's backward-compatibility program (whose
comparison statement is registered as a conjecture with proof required, and
is not advanced by this chapter).

Third, the \emph{ACBE and RPM structures survive as structure}. The descent
maps are still arrows of $\cC$; the section-retraction identity
$\varepsilon_X\circ\eta_X=\id_X$ still holds; the one-directionality of
Remark~\ref{ch05:rem:notadjunction} is untouched. What is withdrawn is one
particular consumption of this structure---as cover generators---not the
structure itself. Indeed the successor re-consumes the same names in two
new, carefully separated roles: as formal coordinate labels on a new finite
context factor (where the mathematics never depends on the labels), and as
cover-kind tags in the external marked-cover records (where provenance
belongs, per Remark~\ref{ch05:rem:design}).

\begin{ruleenv}[Editorial rule derived from the corrigendum]
\label{ch05:rule:editorial}
No artifact of the program may assert or imply that the legacy ACBE or RPM
clauses generate distinct or nontrivial modalities, that the legacy topology
records clause provenance, or that a bare legacy representable is a sheaf.
Permitted statements are: the conditional topos existence for the legacy
site; the equality $J_{\mathrm{legacy}}=J_{\mathrm{js}}$; and the collapse
theorem itself with its registered proof status. Text describing the
\emph{intent} of the legacy clauses must mark it as retracted intent, as
Remark~\ref{ch05:rem:intent} does.
\end{ruleenv}

\section{Executable corroboration}\label{ch05:sec:certificate}

Both collapse mechanisms are small enough to check by machine, and the
registered finite certificate does so: its check
\texttt{legacy\_acbe\_identity\_forces\_maximal\_sieve} constructs a sieve
from a generating family containing an identity and verifies elementwise
that the closure is the maximal sieve, and
\texttt{legacy\_rpm\_split\_epi\_forces\_identity} constructs the sieve
generated by a split epimorphism with an explicit section and verifies that
the identity lands in the closure, hence maximality. These checks live in
the same certificate that exhaustively corroborates the successor's finite
theory, and they are reported as passing in the registered certificate
artifact; the certificate chapter of this monograph documents the full
check roster and its claim boundary. Two remarks keep the corroboration in
proportion. The checks are corroboration of
Theorem~\ref{ch05:thm:collapse}'s mechanisms on executable instances, not a
formal verification of the theorem, whose quantification is over all
registered clause instances; and they are deliberately redundant---the
informal proof is three lines---because the program's evidence discipline
prices redundancy at approximately zero and prices unexecuted claims high.
An auditor who trusts neither the prose nor the author can rerun the
certificate and watch the identity arrive in the closed sieve.

\section{Chapter summary}\label{ch05:sec:summary}

The registered legacy site admits covers from three clauses. The ACBE
clause's reflexivity provision places $\id_X$ among its generators; the RPM
clause's counit is split by its registered section, and closure places
$\varepsilon_X\circ\eta_X=\id_X$ in its sieve. By the two elementary lemmas,
every clause-generated sieve is maximal
(Theorem~\ref{ch05:thm:collapse}(1),(2)); maximal sieves cost nothing, so
$J_{\mathrm{legacy}}=J_{\mathrm{js}}$
(Theorem~\ref{ch05:thm:collapse}(3)); at the presheaf level, every presheaf
satisfies both clauses' gluing conditions automatically, the RPM
amalgamation being manufactured by the section itself
(Proposition~\ref{ch05:prop:autoglue}); and the intended modalities are
trivial (Corollary~\ref{ch05:cor:trivial} via
Proposition~\ref{ch04:prop:trivial}). The collapse is an instance of a
structural fact: topology generation is a closure operator and erases both
clause labels and generator identity (Proposition~\ref{ch05:prop:closure},
Finding~\ref{ch05:fnd:erasure}). The three tempting repairs fail on
schedule: relabeling is label-blind
(Proposition~\ref{ch05:prop:repair1}), sub-sieve surgery is unstable or
vacuous (Proposition~\ref{ch05:prop:repair2},
Remark~\ref{ch05:rem:pretop}), and object adjunction preserves maximality
absolutely (Proposition~\ref{ch05:prop:repair3});
Table~\ref{ch05:tab:repairs} is the ledger. The v7.4 corrigendum retracts
the modal inference, the provenance inference, and silent
representable-as-sheaf usage, while preserving conditional topos existence,
the observational topos, and the ACBE/RPM structure as structure
(Section~\ref{ch05:sec:corrigendum}); Rule~\ref{ch05:rule:editorial}
enforces the boundary in prose, and
Invariant~\ref{ch05:inv:designrule} carries the mathematical lesson
forward: nontrivial covers require proper generators---no identities, no
isomorphisms, no split epimorphisms onto the target.

Two questions remain before the successor construction can begin. The first
is whether the legacy topology, collapsed though it is, at least treats the
legacy objects themselves well---whether representables are sheaves, so
that legacy objects pass into the topos unmodified. Chapter
\ref{ch06:main} answers no: subcanonicity is not derivable from the
registered axioms, by an explicit small countermodel with three objects up
to isomorphism, and the standard
repair routes every legacy object through its associated sheaf. The second
is whether proper generators supporting genuinely distinct, nontrivial,
well-understood modalities can be exhibited at all within the program's
conservative constraints. The remaining chapters answer yes, on a finite
context factor, with the entire finite theory proved directly and
corroborated exhaustively.

\begin{exercise}\label{ch05:exr:variants}
(i) Show that a sieve containing a \emph{monomorphism that is also
covering-split} in the weaker sense that some $h$ with
$\varepsilon\circ h=\id$ exists merely in a larger ambient category
$\cC'\supseteq\cC$ need \emph{not} be maximal in $\cC$, and explain exactly
which step of Lemma~\ref{ch05:lem:splitepi} fails; conclude that the
collapse depends on the section being an arrow of the site itself, and
verify that Axiom~\ref{ch05:ax:R} does place $\eta_X$ in $\cC$.
(ii) Show that on a poset the converse of
Invariant~\ref{ch05:inv:designrule} also holds: a generating family every
member of which is a \emph{proper} subobject of $X$ can never generate the
maximal sieve on $X$, because the generated sieve is
$\bigcup_i\mathord{\downarrow}U_i$ and $X\notin\mathord{\downarrow}U_i$
whenever $U_i\subsetneq X$. Then exhibit a \emph{non}-poset
example---a category carrying a split epimorphism $\varepsilon:E\to X$ that
is not an isomorphism---in which a generating family of non-identity,
non-isomorphism arrows nonetheless generates $m_X$, and identify which
hypothesis of the poset argument fails.
\end{exercise}

\begin{exercise}\label{ch05:exr:closurecount}
On the poset $L_2$ of Chapter~\ref{ch04:main}, compute
$\mathrm{Gen}(G)$ for all $2^{6}$ subsets $G$ of the six sieves on the top
object (holding the null cover fixed), verify idempotency and monotonicity
empirically, and count how many distinct topologies arise. This is the
smallest nontrivial instance of the non-injectivity in
Example~\ref{ch05:ex:noninjective}.
\end{exercise}
